\documentclass[11pt]{article}
\usepackage{rabbithole}

\phasenumber{13}
\phasetitle{The gh CLI in full}
\phasesubtitle{Never leave the terminal, and reach the parts of GitHub that have no command.}

\hypersetup{pdftitle={Rabbit Holes, Phase 13: The gh CLI in full},
            pdfauthor={Accelerate, Manipal University Jaipur},
            pdfsubject={git and GitHub handbook}}

\begin{document}
\maketitlepage

\section{Why bother}

The web interface is fine. The argument for \co{gh} is not speed of a single action,
it is that everything becomes scriptable.

``Add twelve people to a repository'' is twelve pages of clicking, or one loop.
``Which of my pull requests have failing checks'' is a lot of tab switching, or one
command. Phase 3's collaborator invitations were a \co{gh api} call for exactly this
reason.

\section{Auth and config}

\begin{shell}
gh auth login
gh auth status
gh auth token
gh auth refresh --scopes admin:org,workflow
gh auth logout
gh auth switch
\end{shell}

\co{gh auth refresh} is worth knowing. When a command fails with a permissions error,
it is usually a missing scope rather than a real permission problem, and this adds
one without logging out.

\begin{shell}
gh config set editor "code --wait"
gh config set git_protocol ssh
gh config set prompt disabled
gh config list
\end{shell}

\section{Repositories}

\begin{shell}
gh repo create my-project --public --clone
gh repo create --source=. --public --push
gh repo create team-alpha --template accelerate-muj/git-started-team --public --clone
gh repo clone accelerate-muj/git-started
gh repo fork accelerate-muj/git-started --clone
gh repo sync my-username/git-started
gh repo view accelerate-muj/git-started --web
gh repo list accelerate-muj --limit 100
gh repo edit --description "..." --add-topic git --add-topic workshop
gh repo archive old-project
gh repo delete unwanted-project
\end{shell}

\begin{note}
\co{gh repo delete} needs the \co{delete\_repo} scope, which is not granted by
default. That is deliberate friction on an irreversible action, and it is a
reasonable design.
\end{note}

\section{Pull requests}

The largest command group, and the one that saves the most time.

\begin{shell}
gh pr create --fill
gh pr create --title "..." --body "..." --base main --draft
gh pr create --reviewer shashwat-deep --assignee @me --label enhancement
gh pr list --state open --author @me
gh pr list --search "review-requested:@me"
gh pr status
gh pr view 42
gh pr view 42 --comments
gh pr diff 42
gh pr checkout 42
gh pr review 42 --approve
gh pr review 42 --request-changes --body "See the comment on line 40"
gh pr review 42 --comment --body "Two small things"
gh pr merge 42 --squash --delete-branch
gh pr merge 42 --merge --auto
gh pr close 42
gh pr ready 42
gh pr checks 42
gh pr comment 42 --body "Rebased onto main"
\end{shell}

Two of these deserve highlighting.

\co{gh pr checkout 42} fetches the contributor's branch, including from a fork, and
checks it out locally. It handles the remote juggling for you. Reviewing code you can
run is a different activity from reading a diff, and this is what makes it a
one-command activity.

\co{gh pr merge --auto} queues the merge for when checks pass. Enable it and walk
away, rather than watching a CI run for eleven minutes.

\section{Issues, releases, workflows}

\begin{shell}
gh issue create --title "..." --body "..." --label bug --assignee @me
gh issue list --state open --label bug --limit 50
gh issue develop 42 --checkout
\end{shell}

\co{gh issue develop} creates a branch linked to an issue and checks it out. The link
shows on the issue, so the connection exists before the pull request does.

\begin{shell}
gh release create v1.0.0 --generate-notes
gh release upload v1.0.0 ./dist/app.zip
gh release download v1.0.0 --pattern "*.zip"
gh release list
\end{shell}

\begin{shell}
gh workflow list
gh workflow run deploy.yml --ref main -f environment=staging
gh workflow disable deploy.yml
gh run list --workflow=ci.yml --limit 10
gh run view <run-id> --log
gh run view <run-id> --log-failed
gh run watch <run-id>
gh run rerun <run-id> --failed
gh run cancel <run-id>
\end{shell}

\co{gh run view --log-failed} prints only the failing steps' output. On a large
matrix build that is the difference between reading six lines and scrolling through
four thousand.

\section{gh api}

The escape hatch, and the reason \co{gh} can do anything GitHub can.

\begin{shell}
gh api user
gh api repos/accelerate-muj/git-started
gh api repos/:owner/:repo/branches
\end{shell}

\co{:owner} and \co{:repo} are filled in from the current directory's repository.

\subsection{Methods and fields}

\begin{shell}
gh api -X PUT repos/:owner/:repo/collaborators/someone -f permission=push
gh api -X POST repos/:owner/:repo/issues -f title="Bug" -f body="Details"
gh api -X PATCH repos/:owner/:repo -F has_issues=true
gh api -X DELETE repos/:owner/:repo/collaborators/someone
\end{shell}

\begin{trap}
\co{-f} sends the value as a string. \co{-F} lets git interpret it, so booleans and
numbers are sent as booleans and numbers.

\co{-f has\_issues=true} sends the string \co{"true"}, which some endpoints reject
and others silently ignore. Use \co{-F} for anything that is not text. This is a
recurring half hour of confusion for people scripting against the API.
\end{trap}

\subsection{Filtering output}

\co{gh} has \co{jq} built in, so you do not need it installed:

\begin{shell}
gh api repos/:owner/:repo/pulls --jq '.[].title'
gh api repos/:owner/:repo/pulls --jq '.[] | "\(.number)  \(.user.login)  \(.title)"'
gh api repos/accelerate-muj/git-started/contributors --jq '.[].login'
gh api user/repos --paginate --jq '.[] | select(.fork == false) | .name'
\end{shell}

\co{--paginate} follows every page automatically, which matters because the API
returns thirty items by default and quietly truncates otherwise.

\subsection{GraphQL}

Some things, notably Projects, only exist in the GraphQL API:

\begin{shell}
gh api graphql -f query='
  query {
    viewer {
      login
      repositories(first: 5, orderBy: {field: PUSHED_AT, direction: DESC}) {
        nodes { name pushedAt }
      }
    }
  }'
\end{shell}

\begin{tryit}
Phase 3's collaborator loop, for a whole team at once:

\begin{shell}
for user in alice bob charlie dev; do
  gh api -X PUT "repos/$(gh repo view --json nameWithOwner -q .nameWithOwner)/collaborators/$user" -f permission=push
done
\end{shell}

Then check who has accepted:

\begin{shell}
gh api repos/:owner/:repo/invitations --jq '.[].invitee.login'
\end{shell}
\end{tryit}

\section{Aliases}

\begin{shell}
gh alias set prs 'pr list --author @me'
gh alias set review 'pr list --search "review-requested:@me"'
gh alias set co 'pr checkout'
gh alias set bugs 'issue list --label bug --state open'
gh alias list
\end{shell}

Shell aliases with \co{--shell} can take arguments and use pipes:

\begin{shell}
gh alias set --shell teamup 'gh api -X PUT repos/$1/collaborators/$2 -f permission=push'
gh teamup accelerate-muj/git-started someone
\end{shell}

\section{Extensions}

\begin{shell}
gh extension search
gh extension install dlvhdr/gh-dash
gh extension list
gh extension upgrade --all
\end{shell}

\co{gh-dash} is a terminal dashboard of your pull requests and issues, and it is the
one most people keep. Extensions are just repositories named \co{gh-something} with
an executable, so writing one is a short afternoon.

\section{gh inside Actions}

\co{gh} is preinstalled on GitHub-hosted runners. It needs a token in the
environment:

\begin{terminal}
- name: Comment on the PR
  env:
    GH_TOKEN: ${{ secrets.GITHUB_TOKEN }}
  run: gh pr comment "$PR" --body "Build finished"
\end{terminal}

This is exactly how the Phase 2 poem checker works in this repository. Look at
\co{.github/workflows/check-poem.yml} for a complete example.

\begin{trap}
The automatic \co{GITHUB\_TOKEN} is scoped to the repository the workflow runs in,
and its permissions are whatever the \co{permissions:} block grants. If a \co{gh}
call in CI fails with a 403, the fix is almost always adding the right permission to
that block rather than creating a personal token.

Reach for a personal access token only when you genuinely need to touch a
\emph{different} repository, and store it as a secret.
\end{trap}

\section{Reference}

\vspace{0.5em}
\begin{tabularx}{\linewidth}{@{}L{6.6cm} X@{}}
\toprule
\rhtablehead{Goal} & \rhtablehead{Command} \\
\midrule
Repo from the current folder & \co{gh repo create -{}-source=. -{}-public -{}-push} \\
Repo from a template & \co{gh repo create <name> -{}-template <repo>} \\
Check out someone's PR & \co{gh pr checkout 42} \\
Merge when checks pass & \co{gh pr merge 42 -{}-squash -{}-auto} \\
Only the failing CI output & \co{gh run view <id> -{}-log-failed} \\
Branch linked to an issue & \co{gh issue develop 42 -{}-checkout} \\
Add a collaborator & \co{gh api -X PUT repos/:owner/:repo/collaborators/<u>} \\
Every page of a list & \co{gh api <endpoint> -{}-paginate} \\
Boolean field to the API & \co{-F key=true}, never \co{-f} \\
Save a command as a word & \co{gh alias set <name> '<command>'} \\
\bottomrule
\end{tabularx}
\vspace{0.5em}

\checkyourself{
\begin{enumerate}
  \item What does \co{gh pr checkout} do that \co{git fetch} alone does not?
  \item Difference between \co{-f} and \co{-F} on \co{gh api}, and when it bites.
  \item Why do you need \co{-{}-paginate}?
  \item A \co{gh} call in a workflow returns 403. What is the first thing to check?
  \item Write an alias that lists pull requests waiting on your review.
\end{enumerate}}

\vspace{1em}
{\color{rhborder}\rule{\linewidth}{0.8pt}}

\textbf{Next:} Phase 14, \emph{Actions and CI/CD}, on making the repository do work
by itself: running tests on every pull request, and deploying when you tag a release.

\end{document}
